\documentclass[conference]{IEEEtran}
\usepackage[utf8]{inputenc}
\usepackage[backend=biber, style=authoryear]{biblatex}
\usepackage[T1]{fontenc}
\usepackage[brazil]{babel}
\usepackage{amsmath}
\usepackage{graphicx}
\usepackage{booktabs}
\usepackage{float}
\usepackage{enumitem}
\usepackage[hidelinks]{hyperref}
\addbibresource{references.bib}

% Caixa provisória para figuras ainda não geradas.
% Uso: \placeholder{descrição}  ou  \placeholder[\textwidth]{descrição}
\newcommand{\placeholder}[2][\columnwidth]{%
	\fbox{\parbox[c][3cm][c]{\dimexpr#1-2\fboxsep-2\fboxrule\relax}{%
			\centering\footnotesize\textit{#2}}}}

% Tabelas e figuras finais, geradas por `make analise` em results/final/.
% Figura: \figuraresultado[largura]{arquivo.png}{descrição}; sem o arquivo, usa o placeholder.
\newcommand{\figuraresultado}[3][\columnwidth]{%
	\IfFileExists{results/final/#2}%
	{\includegraphics[width=#1]{results/final/#2}}%
	{\placeholder[#1]{#3}}}
% Linhas de uma tabela, entre \midrule e \bottomrule: \linhasresultado{tab-principal}.
% Usa o \input primitivo, que funciona dentro de um tabular.
\makeatletter
\newcommand{\linhasresultado}[1]{\@@input results/final/#1.tex }
\makeatother

% Convenções:
%   - rótulos com prefixo: sec:, subsec:, fig:, tab:, eq:
%   - tabelas e figuras geradas por código ficam em results/final/ (ver README,
%     "Reprodução"); os CSVs brutos ficam em results/raw/
%   - orçamento total sugerido pelo enunciado: ~5 páginas sem anexos;
%     figuras extras vão para anexos.tex

\title{
	Busca Local e Metaheurísticas no Sequenciamento de Manutenção de Ativos \\
	\textit{\large TP2 - Fundamentos de Inteligência Artificial (2026/2)}
}
\author{
	\IEEEauthorblockN{Diogo Muzzi Mortimer}
	\IEEEauthorblockA{\textit{Universidade Federal de Minas Gerais}}
	\IEEEauthorblockA{2022055394}
	\and
	\IEEEauthorblockN{Paulo Henrique Soares Gomes}
	\IEEEauthorblockA{\textit{Universidade Federal de Minas Gerais}}
	\IEEEauthorblockA{2024433370}
	\and
	\IEEEauthorblockN{Pedro Pessoa Santos Brochado}
	\IEEEauthorblockA{\textit{Universidade Federal de Minas Gerais}}
	\IEEEauthorblockA{2023093109}
}

% Conteudo
\begin{document}

\maketitle

\begin{abstract}
Este trabalho avalia o desempenho de heurísticas de Busca Local (Greedy Descent com Random Restart e Simulated Annealing) e de uma metaheurística populacional (Algoritmo Genético) na resolução do problema de sequenciamento de manutenção preventiva de ativos. O objetivo consiste em minimizar o atraso ponderado total $h(\pi)$ da operação. A avaliação foi conduzida sob um orçamento computacional fixo de avaliações da função objetivo, utilizando três instâncias de diferentes magnitudes ($n=12$, $50$ e $150$), com 30, 30 e 10 execuções independentes, respectivamente. A análise estatística foi fundamentada no teste não paramétrico de Mann-Whitney. Os resultados revelam que, embora as abordagens de busca local apresentem eficiência notável e baixo custo por iteração nas instâncias menores, o Algoritmo Genético demonstra superioridade estatística significativa na instância de grande escala, justificando seu maior custo computacional ao manter um equilíbrio otimizado entre exploração e intensificação no vasto espaço de busca.
\end{abstract}

\begin{IEEEkeywords}
Busca local, simulated annealing, algoritmo genético, sequenciamento, otimização combinatória
\end{IEEEkeywords}
% Enunciado: Seções 1.1 e 1.2 (o problema é dado; apresentar de forma breve)
% Tamanho sugerido: ~2/3 de coluna

\section{Introdução}
\label{sec:introducao}

% Parágrafo 1 -- contexto: concessionária de energia, uma única equipe que
% parte da base e atende um ativo por vez; cada ativo j tem posição
% (x_j, y_j), duração p_j, prazo d_j e criticidade w_j.

% Formulação compacta, definindo a notação usada no resto do texto:
% - tempo de deslocamento t_ij arredondado (\label{eq:deslocamento});
% - instante de conclusão C_{pi_k} (\label{eq:conclusao});
% - função de avaliação h(pi) = sum w_j max(0, C_j - d_j) (\label{eq:fitness}).
% Uma frase sobre o tamanho do espaço (n!, n = 150 => >260 algarismos), o
% fato de toda permutação ser viável e a tensão entre ir ao mais próximo e
% ir ao mais urgente.

% Parágrafo 2 -- objetivo: comparar busca local (greedy descent com random
% restart, simulated annealing) e uma metaheurística populacional
% (algoritmo genético) sob o mesmo orçamento de avaliações, observando o
% equilíbrio exploração x intensificação. Citar \cite{poole2010}.

% Parágrafo 3 -- organização do texto, referenciando:
% Seção~\ref{sec:metodologia}, Seção~\ref{sec:algoritmos},
% Seção~\ref{sec:resultados} (parte central) e Seção~\ref{sec:conclusao}.

\section{Metodologia}
\label{sec:metodologia}

Para assegurar rigor e paridade experimental, a arquitetura de \textit{software} implementada força os três algoritmos estruturados a interagirem com o domínio do problema unicamente através de uma interface genérica de métodos, exposta na Tabela \ref{tab:interface}. O consumo computacional é estritamente controlado por um orçamento estático $B$ de avaliações da função objetivo $h(\pi)$. Essa validação é instrumentada em um único ponto base da arquitetura, que aciona o encerramento do processamento para qualquer método que atinja seu respectivo limite.

\begin{table}[H]
\centering
\caption{Interface comum entre problema e algoritmos.}
\label{tab:interface}
\begin{tabular}{ll}
\toprule
Método                          & Papel \\
\midrule
\texttt{solucao\_aleatoria()}   & Gera um ponto de partida inicial \\
\texttt{h(s)}                   & Função de avaliação (fitness) \\
\texttt{sucessores(s)}          & Itera sobre os adjacentes do espaço de busca \\
\texttt{sucessor\_aleatorio(s)} & Sorteia um vizinho único da permutação \\
\texttt{cruzamento(s1, s2)}     & Operador de recombinação de pais \\
\texttt{mutacao(s)}             & Operador de perturbação estocástica \\
\bottomrule
\end{tabular}
\end{table}

A progressão e geração de sucessores no espaço de busca local ocorrem por meio de três operadores clássicos de vizinhança: \textit{Swap} (troca de posição direta entre dois ativos iterados), \textit{Insertion} (remoção de um ativo com realocação em índice distinto) e \textit{2-opt} (inversão integral do sentido de um subtrecho). A Tabela \ref{tab:tamanho-vizinhancas} define o quantitativo da dimensão da vizinhança $|N(\pi)|$ gerada analiticamente por cada mecanismo. A avaliação destas abordagens ditou o operador eleito nos ensaios finais, discutido na Seção \ref{subsec:busca-local}.

\begin{table}[H]
\centering
\caption{Tamanho da vizinhança induzida por cada operador.}
\label{tab:tamanho-vizinhancas}
\begin{tabular}{lcccc}
\toprule
Operador  & $|N(\pi)|$ & $n=12$ & $n=50$ & $n=150$ \\
\midrule
Swap      & $n(n-1)/2$ & 66 & 1.225 & 11.175 \\
Insertion & $(n-1)^2$ & 121 & 2.401 & 22.201 \\
2-opt     & $n(n-1)/2$ & 66 & 1.225 & 11.175 \\
\bottomrule
\end{tabular}
\end{table}

Toda a infraestrutura da função de avaliação foi validada unitariamente em teste algorítmico, asseverando, por exemplo, que a permutação estritamente identidade do cenário mínimo ($n=12$) resulte de forma precisa na aptidão base $h = 4366$.

A formulação obedece ao arranjo iterativo e orçamento de avaliações listados na Tabela \ref{tab:instancias}. Como suporte referencial para as discussões textuais subsequentes (onde as afirmações exigem consistência científica), os cenários executam rodadas independentes para isolar distorções aleatórias. Tais afirmações no texto estarão invariavelmente baseadas no teste estatístico não paramétrico de Mann-Whitney (nível de significância base), acompanhados diretamente pelo relato pontual de seus respectivos \textit{p-valores}.

\begin{table}[H]
\centering
\caption{Instâncias, orçamento e número de execuções independentes.}
\label{tab:instancias}
\begin{tabular}{lccccc}
\toprule
Instância     & $n$ & $\sum p_j$ & $h_{\text{ref}}$ & $B$ & Execuções \\
\midrule
pequena       & 12  & 55         & 1.875            & 50.000  & 30        \\
intermediária & 50  & 239        & 16.332           & 150.000 & 30        \\
grande        & 150 & 684        & 60.417           & 300.000 & 10        \\
\bottomrule
\end{tabular}
\end{table}
% Enunciado: Seções 2.3--2.5, avaliação incremental (Seção 3) e T2
% Tamanho sugerido: ~1,5 coluna no total. A parte central do relatório é a
% análise comparativa: aqui só o necessário para entender e reproduzir cada
% algoritmo; resultados de calibração e diagnóstico em poucas linhas, com
% material extra em anexos.tex.

\section{Algoritmos}
\label{sec:algoritmos}

% Uma ou duas frases: os três algoritmos usam a interface da
% Tabela~\ref{tab:interface} e o mesmo contador de avaliações.

\subsection{Greedy Descent com Random Restart}
\label{subsec:busca-local}

O método de descida gulosa (\textit{Greedy Descent}) foi implementado com reinícios aleatórios (\textit{Random Restart}) para evitar que o algoritmo fique preso em ótimos locais. A partir de uma permutação aleatória inicial, o algoritmo explora a vizinhança e sempre se move para o sucessor de menor custo $h$. Quando chega a um ótimo local (onde nenhum vizinho melhora o custo), a execução é reiniciada a partir de uma nova solução gerada do zero. Esse ciclo guarda o registro da melhor solução global encontrada e continua até que o orçamento de avaliações $B$ acabe.

Implementamos três variantes para escolher o próximo passo: a) \textit{Best improvement}, que varre toda a vizinhança e escolhe o melhor sucessor; b) \textit{First improvement}, que aceita o primeiro vizinho que melhore o custo $h$; e c) \textit{Sucessor aleatório com aceitação de melhora}, que sorteia apenas um vizinho e avança se houver melhora.

Para otimizar o tempo de execução, implementamos a avaliação incremental. Como os movimentos de vizinhança não alteram o prefixo da permutação, o instante de conclusão $C_k$ e o custo acumulado desse trecho inicial não mudam. O nosso código reaproveita esses valores e recalcula o custo apenas do sufixo alterado. Os ganhos de velocidade e a validação exata dos cálculos estão detalhados na Tabela \ref{tab:incremental}.

\begin{table}[H]
\centering
\caption{Tempo médio por avaliação: completa x incremental.}
\label{tab:incremental}
\begin{tabular}{llccc}
\toprule
Instância & Operador & Completa ($\mu$s) & Incremental ($\mu$s) & Speedup \\
\midrule
pequena & Swap & 1.14 & 1.04 & 1.10 \\
 & Insertion & 1.06 & 1.04 & 1.02 \\
 & 2-opt & 1.11 & 0.98 & 1.14 \\
\midrule
intermedi\'aria & Swap & 3.54 & 2.83 & 1.25 \\
 & Insertion & 3.53 & 2.81 & 1.26 \\
 & 2-opt & 3.58 & 2.79 & 1.28 \\
\midrule
grande & Swap & 10.23 & 7.28 & 1.40 \\
 & Insertion & 10.21 & 7.29 & 1.40 \\
 & 2-opt & 10.13 & 7.36 & 1.38 \\
\bottomrule
\end{tabular}
\end{table}

A qualidade geométrica das vizinhanças foi testada sob o mesmo orçamento, mostrando qual operador encontra os melhores ótimos locais (Tabela \ref{tab:vizinhancas}). Baseado nisso, adotamos o \textit{Swap} para os testes finais.

\begin{table}[H]
\centering
\caption{Greedy descent por vizinhança: melhor $h$ (média $\pm$ dp / mínimo).}
\label{tab:vizinhancas}
\begin{tabular}{lccc}
\toprule
Operador  & pequena & intermediária & grande \\
\midrule
Swap & 1875.0 $\pm$ 0.0/1875 & 8575.4 $\pm$ 3318.2/21887 & 390602.0 $\pm$ 46825.8 \\
Insertion & 1877.1 $\pm$ 11.3/1875 & 21313.6 $\pm$ 2867.6/16395 & 697380.3 $\pm$ 55287.4 \\
2-opt & 1875.0 $\pm$ 0.0/1875 & 20323.7 $\pm$ 1323.4/18026 & 522030.8 $\pm$ 34483.7 \\
\bottomrule
\end{tabular}
\end{table}

A comparação entre as três abordagens de melhoria sob um orçamento fixo mostra o problema do custo computacional: na instância grande ($n=150$), fazer uma varredura completa da vizinhança \textit{Swap} gasta 11.175 avaliações, o que limita muito a quantidade de passos que o algoritmo consegue dar (Tabela \ref{tab:gd-passos}). 

\begin{table}[H]
\centering
\caption{Passos de melhoria e restarts por variante (média por execução).}
\label{tab:gd-passos}
\begin{tabular}{llccc}
\toprule
Instância & Variante & Passos & Restarts & Melhor $h$ \\
\midrule
pequena & best & 667.1 & 88.9 & 1875.0 \\
 & first & 1718.5 & 88.4 & 1875.0 \\
 & aleat\'oria & 3691.0 & 255.3 & 1877.1 \\
\midrule
intermedi\'aria & best & 120.0 & 2.0 & 29388.6 \\
 & first & 346.6 & 1.2 & 28575.4 \\
 & aleat\'oria & 2460.9 & 19.2 & 24690.9 \\
\midrule
grande & best & 26.0 & 0.0 & 477590.1 \\
 & first & 417.8 & 0.0 & 390602.0 \\
 & aleat\'oria & 1720.4 & 2.3 & 156607.1 \\
\bottomrule
\end{tabular}
\end{table}
\subsection{Simulated Annealing}
\label{subsec:tempera}

O algoritmo de \textit{Simulated Annealing} (SA) foi implementado sorteando vizinhos através do método \texttt{sucessor\_aleatorio(s)}. Se a nova solução diminuir ou mantiver o atraso ponderado atual, o movimento é aceito. Se a nova solução for pior ($\Delta > 0$), a aceitação ocorre com base na probabilidade da distribuição de Gibbs, calculada por $e^{-\Delta/T}$.

A calibração do resfriamento (\textit{annealing schedule}) é feita automaticamente para cada instância. Antes da busca principal começar, o algoritmo faz um passeio aleatório de 200 passos (que são descontados do orçamento) para medir o aumento médio de custo $\overline{\Delta^+}$ daquela instância. Com base nisso, a temperatura inicial $T_0$ é calculada para garantir 90\% de aceitação de pioras no início ($p_0 = 90\%$), e a final $T_f$ para ter apenas 0,1\% de aceitação no fim ($p_f = 0,1\%$). O coeficiente de resfriamento geométrico $\alpha$ ($T \leftarrow \alpha T$) é calculado para que a temperatura atinja $T_f$ exatamente quando o orçamento de avaliações acabar (Tabela \ref{tab:sa-parametros}).

\begin{table}[H]
\centering
\caption{Parâmetros do simulated annealing por instância.}
\label{tab:sa-parametros}
\begin{tabular}{lccccc}
\toprule
Instância     & $\overline{\Delta^+}$ & $T_0$ & $T_f$ & $\alpha$ & Iter./$T$ \\
\midrule
pequena & 770.4 & 7311.8 & 111.52 & 0.999916 & 1 \\
intermedi\'aria & 4023.9 & 38191.9 & 582.52 & 0.999972 & 1 \\
grande & 10465.2 & 99327.2 & 1514.99 & 0.999986 & 1 \\
\bottomrule
\end{tabular}
\end{table}

A calibração individual por tamanho do problema é fundamental, pois a ordem de grandeza das multas cresce rapidamente conforme $n$ aumenta. A Figura \ref{fig:sa-aceitacao} mostra que o SA começa praticamente como um \textit{random walk} (aceitando muitas pioras) e termina como uma descida gulosa (quase não aceitando pioras), provando que a calibração funcionou corretamente.

\begin{figure}[H]
\centering
\figuraresultado{sa_aceitacao.png}{Taxa de aceitação de pioras (janela móvel) x avaliação (mediana das execuções), uma curva por instância; T em eixo secundário opcional.}
\caption{Taxa de aceitação de movimentos piores ao longo da execução do simulated annealing.}
\label{fig:sa-aceitacao}
\end{figure}
\subsection{Algoritmo Genético}
\label{subsec:genetico}

No Algoritmo Genético implementado, cada indivíduo representa uma rota completa de manutenção válida. Cada gene é a posição no agendamento, e o alelo é o identificador do ativo a ser visitado. Como a função $h$ calcula multas financeiras (logo, o objetivo é minimizar), utilizamos a seleção por torneio: sorteamos pequenos grupos da população e o indivíduo com o menor custo é escolhido para se reproduzir.

Usar a codificação por permutação impede o uso de cruzamentos simples (de 1 ou $k$ pontos), pois isso faria com que ativos fossem duplicados ou esquecidos. Para resolver isso, usamos o operador \textit{Order Crossover} (OX), que preserva a validade das rotas. Ele copia um trecho contínuo do primeiro pai e preenche o restante do cromossomo respeitando a ordem circular dos ativos do segundo pai. A Tabela \ref{tab:ag-cruzamento} mostra um exemplo passo a passo.

\begin{table}[H]
\centering
\caption{Exemplo passo a passo do cruzamento \textit{Order Crossover} (OX) adotado.}
\label{tab:ag-cruzamento}
\begin{tabular}{ll}
\toprule
Etapa                   & Permutação \\
\midrule
Pai 1                   & (3, \textbf{1, 5, 2}, 4) \\
Pai 2                   & (4, 2, 1, 5, 3) \\
Segmento herdado        & (--, \textbf{1, 5, 2}, --) \\
Preenchimento (Pai 2)   & (4, --, --, --, 3) \\
Filho                   & (4, 1, 5, 2, 3) \\
\bottomrule
\end{tabular}
\end{table}

\begin{table}[H]
\centering
\caption{Parâmetros do algoritmo genético por instância.}
\label{tab:ag-parametros}
\begin{tabular}{lccccc}
\toprule
Instância     & $N$ & $P_c$ & $P_m$ & Torneio & Elite \\
\midrule
pequena & 100 & 0.90 & 0.20 & 3 & 2 \\
intermedi\'aria & 100 & 0.90 & 0.20 & 3 & 2 \\
grande & 100 & 0.90 & 0.20 & 3 & 2 \\
\bottomrule
\end{tabular}
\end{table}

A mutação é feita aplicando o operador \textit{Swap} com uma probabilidade baixa ($P_m$). O elitismo também foi implementado para garantir que as 2 melhores rotas da população passem diretamente para a geração seguinte sem sofrer alterações. O algoritmo para imediatamente ao atingir o limite orçamental $B$, cortando gerações pela metade se necessário.

Monitorar a diversidade da população (número de cromossomos distintos, mostrado na Figura \ref{fig:ag-diversidade}) é crucial: se a população convergir rápido demais para soluções iguais, o algoritmo perde a capacidade de busca e vira apenas uma busca local cara.

\begin{figure}[H]
\centering
\figuraresultado{ag_diversidade.png}{Diversidade da população x geração (mediana das execuções), uma curva por instância.}
\caption{Evolução da diversidade da população do algoritmo genético.}
\label{fig:ag-diversidade}
\end{figure}

% Enunciado: Seções 3 (estatística, curvas) e 5
% Tamanho sugerido: ~2 páginas no total (parte central do relatório)

\section{Resultados e Discussão}
\label{sec:resultados}

% Enunciado: Seção 5, itens 1--5
% Tamanho sugerido: ~1 página (figuras de largura dupla)

\subsection{Comparação dos algoritmos}
\label{subsec:comparacao}

% Qualidade das soluções: apresentar a Tabela~\ref{tab:principal} e a
% Figura~\ref{fig:boxplots}.

\begin{table*}[t]
	\centering
	\caption{Melhor $h$ por combinação algoritmo $\times$ instância. Gap em
		relação a $h^\star = 1875$ apenas na instância pequena.}
	\label{tab:principal}
	\begin{tabular}{llccccc}
		\toprule
		Instância     & Algoritmo & Média $\pm$ dp & Mediana & Mínimo & Gap médio & Gap mínimo \\
		\midrule
		pequena       & GD+RR     & --             & --      & --     & --        & --         \\
		              & SA        & --             & --      & --     & --        & --         \\
		              & AG        & --             & --      & --     & --        & --         \\
		\midrule
		intermediária & GD+RR     & --             & --      & --     & \multicolumn{2}{c}{--} \\
		              & SA        & --             & --      & --     & \multicolumn{2}{c}{--} \\
		              & AG        & --             & --      & --     & \multicolumn{2}{c}{--} \\
		\midrule
		grande        & GD+RR     & --             & --      & --     & \multicolumn{2}{c}{--} \\
		              & SA        & --             & --      & --     & \multicolumn{2}{c}{--} \\
		              & AG        & --             & --      & --     & \multicolumn{2}{c}{--} \\
		\bottomrule
	\end{tabular}
\end{table*}

\begin{figure*}[t]
	\centering
	\placeholder[\textwidth]{Boxplots do melhor $h$: três painéis (pequena,
		intermediária, grande), três caixas por painel (GD+RR, SA, AG); linha
		tracejada com o valor de referência. Arquivo: results/boxplots.png}
	\caption{Distribuição do melhor $h$ por algoritmo em cada instância.}
	\label{fig:boxplots}
\end{figure*}

% Testes de hipótese: um p-valor para cada comparação par a par afirmada no
% texto (Tabela~\ref{tab:pvalores}).

\begin{table}[H]
	\centering
	\caption{Mann--Whitney par a par (p-valor; algoritmo melhor entre parênteses).}
	\label{tab:pvalores}
	\begin{tabular}{lccc}
		\toprule
		Instância     & GD $\times$ SA & GD $\times$ AG & SA $\times$ AG \\
		\midrule
		pequena       & --             & --             & --             \\
		intermediária & --             & --             & --             \\
		grande        & --             & --             & --             \\
		\bottomrule
	\end{tabular}
\end{table}

% Convergência: descrever as curvas da Figura~\ref{fig:convergencia}; a
% leitura da distância melhor x média do AG entra nas subseções seguintes.

\begin{figure*}[t]
	\centering
	\placeholder[\textwidth]{Curvas de convergência: três painéis (um por
		instância); eixo x = avaliações, eixo y = melhor $h$ (mediana das
		execuções); uma curva por algoritmo + curva da média da população do
		AG. Arquivo: results/convergencia.png}
	\caption{Melhor $h$ em função do número de avaliações.}
	\label{fig:convergencia}
\end{figure*}

% Tempo de execução: mesma contagem de avaliações, tempos de parede
% diferentes (Tabela~\ref{tab:tempo}); de onde vem a diferença (avaliação
% incremental, overhead de população etc.).

\begin{table}[H]
	\centering
	\caption{Tempo médio de parede por execução e avaliações.}
	\label{tab:tempo}
	\begin{tabular}{llccc}
		\toprule
		Instância & Algoritmo & Avaliações & Tempo (s) & Aval./s \\
		\midrule
		--        & --        & --         & --        & --      \\
		\bottomrule
	\end{tabular}
\end{table}

% Enunciado: Seção 5, pontos obrigatórios da discussão
% Tamanho sugerido: ~1 coluna

\subsection{Ótimos locais}
\label{subsec:otimos-locais}
% - Quantos ótimos locais distintos o greedy descent encontrou na
%   instância pequena ao longo dos restarts (somando as 30 execuções).
% - Em quantas execuções cada algoritmo alcançou h* = 1875.
% - Relacionar com ótimo local x ótimo global visto em aula.
% - Opcional (anexo): histograma dos valores de h dos ótimos locais.

\begin{table}[H]
	\centering
	\caption{Instância pequena: ótimos locais e execuções que atingem $h^\star$.}
	\label{tab:otimos-locais}
	\begin{tabular}{lcc}
		\toprule
		Algoritmo & Ótimos locais distintos & Exec. com $h^\star$ (de 30) \\
		\midrule
		\linhasresultado{tab-otimos-locais}
		\bottomrule
	\end{tabular}
\end{table}

\subsection{Exploração $\times$ intensificação}
\label{subsec:exploracao}
% - Para cada algoritmo, qual mecanismo explora e qual intensifica.
% - Onde os dados mostram equilíbrio bem/mal ajustado, apontando figuras:
%   restarts (Tabela~\ref{tab:gd-passos}), taxa de aceitação
%   (Figura~\ref{fig:sa-aceitacao}), diversidade
%   (Figura~\ref{fig:ag-diversidade}) e distância melhor x média do AG
%   (Figura~\ref{fig:convergencia}).

\begin{table}[H]
	\centering
	\caption{Mecanismos de exploração e intensificação de cada algoritmo.}
	\label{tab:exploracao}
	\begin{tabular}{lp{0.33\columnwidth}p{0.33\columnwidth}}
		\toprule
		Algoritmo & Exploração & Intensificação \\
		\midrule
		GD+RR     & --         & --             \\
		SA        & --         & --             \\
		AG        & --         & --             \\
		\bottomrule
	\end{tabular}
\end{table}

\subsection{Efeito da escala}
\label{subsec:escala}
% - O ranking dos três algoritmos muda da instância pequena para a grande?
% - Se sim: hipótese ligada ao custo de varrer a vizinhança (O(n^2)
%   sucessores por passo) e ao tamanho do espaço de busca, sustentada com
%   os dados (Tabelas~\ref{tab:principal}, \ref{tab:gd-passos} e
%   \ref{tab:tempo}).
% - Se não: mostrar com os p-valores que a estabilidade do ranking é
%   estatisticamente sustentada.
% Sem figura própria; reaproveitar as da Seção~\ref{subsec:comparacao}.

\subsection{Quando o algoritmo genético compensa?}
\label{subsec:quando-ag}
% - Existe instância em que o AG supera as duas buscas locais com
%   significância estatística? Citar Tabela~\ref{tab:pvalores}.
% - Hipótese em termos da codificação: a recombinação consegue combinar bons
%   trechos de sequência de dois pais num filho melhor?
% - Experimento adicional que testa a hipótese: desenho (ex.: AG com vs.
%   sem cruzamento, OX vs. PMX, ou taxa de filhos melhores que ambos os
%   pais ao longo das gerações), resultado e p-valor.

\begin{table}[H]
	\centering
	\caption{Experimento adicional: efeito da recombinação.}
	\label{tab:experimento-ag}
	\begin{tabular}{llcc}
		\toprule
		Instância & Configuração & Melhor $h$ (média $\pm$ dp) & p-valor \\
		\midrule
		\linhasresultado{tab-experimento-ag}
		\bottomrule
	\end{tabular}
\end{table}

% Opcional, se o experimento for sobre a dinâmica ao longo das gerações:
% figura com a curva medida (ex.: fração de filhos melhores que os pais
% x geração). Arquivo sugerido: results/ag\_experimento.png


% Tamanho sugerido: ~1/4 de coluna

\section{Conclusão}
\label{sec:conclusao}

% Parágrafo 1 -- retomar o objetivo e responder em poucas frases: qual
% algoritmo foi melhor em cada escala, quão perto da referência chegamos e
% quando o AG compensa. Sem informação nova.

% Parágrafo 2 -- limitações (ex.: calibração limitada, poucas execuções na
% instância grande) e extensões possíveis.


\printbibliography

% Anexos: material que não cabe nas ~5 páginas. Não conta no limite.
% Apague um bloco que ficar vazio.

\appendices

\section{Busca local: material complementar}
\label{anexo:busca-local}
% Ex.: curvas de convergência por vizinhança (T2), histograma dos ótimos
% locais da instância pequena.

\section{Simulated annealing: material complementar}
\label{anexo:tempera}
% Ex.: taxa de aceitação de uma têmpera mal calibrada (T alta demais / baixa
% demais) para contraste com a Figura~\ref{fig:sa-aceitacao}.

\section{Algoritmo genético: material complementar}
\label{anexo:genetico}
% Ex.: diversidade por instância em painéis separados, sensibilidade a N e Pm.


\end{document}
